\chapter{Formato dell'esame e regole del corso}
\label{cap:formato}

\begin{nota}[Fonti]
Slide \emph{0-Intro-Course}, quanto detto in aula alla \rif{2026-09-21}{21/09} (presentazione del
corso) e alla \rif{2026-10-20}{20/10} (linee guida per il report), testi dei laboratori pubblicati
sul portale (sezione ``Lab Rules''). Dove le fonti non coincidono lo segnaliamo nei riquadri
``Da verificare'': controllate sul portale della didattica o chiedete al docente prima dell'esame.
\end{nota}

\begin{verifica}
\begin{itemize}
  \item \textbf{Punteggio dello scritto}: la slide dice \textbf{32 punti}; in aula il docente ha
        detto \textbf{33}.
  \item \textbf{Dimensione dei gruppi}: in aula, in una slide (``Classes' structure'') e nei testi dei
        lab \textbf{4 persone} (``novità di quest'anno''); un'altra slide (``Lab activities'') dice
        \textbf{3 persone}.
  \item \textbf{Scadenza dei report}: la slide e il testo del Lab~3 dicono \textbf{10 giorni} prima
        dell'appello; il testo del Lab~1 dice ``fino a \textbf{due settimane} prima'' di un appello.
  \item \textbf{Lab 2}: la slide lo chiama \en{supervised traffic classification}; in aula è stato
        descritto come classificazione malware/goodware a partire da sequenze di API call (vedi
        cap.~\ref{cap:lab2}).
\end{itemize}
La guida del corso sul portale (\code{didattica.polito.it}) richiede l'accesso autenticato: non è
stato possibile consultarla per risolvere i dubbi.
\end{verifica}

\section{Dati generali}

\begin{center}
\begin{tabularx}{\textwidth}{>{\bfseries}l X}
\toprule
Insegnamento & AI and Cybersecurity (codici 01GYZUV, 01GYZUW) \\
CFU & 6 CFU $\approx$ 60 ore \\
Ripartizione ore & 39 h di lezione (teoria + esempi di codice) + 21 h di laboratorio \\
Docenti & Prof.\ \textbf{Danilo Giordano} (titolare, gran parte del corso) --- \code{danilo.giordano@polito.it} \\
        & Prof.\ \textbf{Matteo Boffa} (NLP e ultima parte del corso) --- \code{matteo.boffa@polito.it} \\
Lingua & Inglese (lezioni e slide) \\
Materiale & Portale della didattica: slide, notebook Python, testi dei lab, registrazioni \\
Ambiente & Python + PyTorch; consigliato \textbf{Google Colab} (GPU gratuita a tempo limitato,
           collaborazione sul notebook). Alternative: Jupyter locale / Anaconda \\
\bottomrule
\end{tabularx}
\end{center}

\subsection{Orari e aule}
\begin{center}
\begin{tabular}{lll}
\toprule
\textbf{Giorno} & \textbf{Orario} & \textbf{Aula / note}\\
\midrule
Lunedì & 13:00--14:30 (1,5 h) & 2C --- lezione \\
Martedì & 16:00--17:30 (1,5 h) & 8I --- \textbf{opzionale}: si usa solo se annunciato sul portale \\
Mercoledì & 11:30--14:30 (3 h) & 2S --- lezione o, quando possibile, laboratorio \\
\bottomrule
\end{tabular}
\end{center}

\subsection{Materiale e libri}
\begin{itemize}
  \item Le \textbf{slide} sono volutamente verbose e \textbf{autosufficienti}: secondo il docente
        bastano per prendere anche il voto massimo. Vengono caricate a ridosso della lezione.
  \item Le lezioni sono registrate; dei laboratori si registra solo la spiegazione iniziale.
  \item Libri consigliati (non necessari):
  \begin{itemize}
    \item I.\,H.\ Sarker, \emph{AI-Driven Cybersecurity and Threat Intelligence}, Springer, 2024.
    \item S.\,J.\,D.\ Prince, \emph{Understanding Deep Learning}, MIT Press, 2023 (molte figure
          delle slide vengono da qui).
  \end{itemize}
\end{itemize}

\subsection{Prerequisiti}
Concetti base di data science e machine learning (il corso fa un rapido ripasso iniziale);
Python con \code{numpy}, \code{pandas}, \code{sklearn}, \code{matplotlib}. PyTorch viene
introdotto durante il corso. Corso utile come base (non obbligatorio):
\emph{Machine Learning for Networking} (01DSMBG, Prof.\ Luca Vassio).

\subsection{Programma}
\begin{multicols}{2}
\textbf{Lezioni (39 h)}
\begin{itemize}
  \item Introduzione ad AI e ML per la cybersecurity
  \item Fondamenti di deep learning
  \item Architetture deep: teoria e implementazione (FFNN, RNN, CNN, GNN)
  \item Supervised e unsupervised, anomaly detection
  \item Introduzione al deep NLP e alle tecniche di embedding
  \item Casi d'uso di cybersecurity (seminari)
  \item Simulazione d'esame (ultima settimana)
\end{itemize}
\columnbreak
\textbf{Laboratori (21 h)}
\begin{enumerate}
  \item Introduzione ai framework di deep learning (3 h)
  \item Classificazione supervisionata (6 h)
  \item Unsupervised e anomaly detection, es.\ rilevamento DDoS (6 h)
  \item NLP e clustering per il rilevamento di botnet (6 h)
\end{enumerate}
\end{multicols}

\section{Laboratori e report}

\subsection{Gruppi}
\begin{itemize}
  \item Gruppi \textbf{auto-organizzati} di \textbf{4 persone} (vedi il riquadro ``Da verificare'':
        una slide dice 3), da registrare con il form Moodle entro fine ottobre.
  \item Il tempo in aula non basta per finire un lab: si completa a casa.
\end{itemize}

\subsection{Cosa consegnare (per ogni lab, Lab 1 incluso)}
Un \textbf{file zip} per lab, chiamato \code{groupID\_labX.zip} secondo i testi dei lab, con:
\begin{enumerate}
  \item \textbf{Report} di \textbf{massimo 10 pagine} (\code{groupID\_labX.pdf}): cosa e come avete
        fatto, decisioni, esperimenti e risultati, con figure e tabelle; frammenti di codice solo se
        servono. Nell'intestazione: \textbf{ID del lab, ID del gruppo, nomi ed email} dei membri.
  \item \textbf{Notebook Jupyter} (nessun limite di pagine), \textbf{eseguito} (output e numeri di
        esecuzione visibili), \textbf{commentato} e \textbf{rieseguibile} (seed fissati); secondo i testi
        dei lab va incluso sia il \code{.ipynb} sia un export HTML.
  \item \textbf{Experiment logbook} (novità di quest'anno): una tabella con una riga per ogni
        addestramento, con budget di circa 25 run per lab; \textbf{non conta} nelle 10 pagine
        (dettagli nel cap.~\ref{cap:lab-report}).
\end{enumerate}

\subsection{Scadenze e valutazione}
\begin{itemize}
  \item \textbf{Lab 1} (solo formativo): consegna \textbf{obbligatoria} per far valutare gli altri lab;
        è valutato solo valido/non valido. Se si consegna una soluzione ``ragionevole''
        \textbf{entro il 1° dicembre} si ottiene \textbf{+1 punto bonus} sul voto finale.
        Se va riconsegnato, o si consegna dopo il 1° dicembre, il bonus si perde.
  \item \textbf{Altri lab}: scadenza \textbf{10 giorni prima della data di un appello} (es.\ appello
        il 15 $\Rightarrow$ consegna entro il 5); il testo del Lab~1 parla invece di due settimane
        (vedi riquadro iniziale).
  \item Si valuta \textbf{il report}; il notebook si guarda solo in caso di dubbi. Il report deve
        quindi essere \textbf{autocontenuto}.
  \item Ogni report vale \textbf{da 0 a 32}; \textbf{tutti} devono avere voto $\geq 18$.
        \begin{itemize}
          \item voto $\geq 18$: \textbf{non} si può riconsegnare per migliorarlo;
          \item voto $< 18$: si \textbf{deve} riconsegnare.
        \end{itemize}
  \item I voti dei report valgono \textbf{2 anni accademici} (fino all'ultimo appello 2027/2028).
  \item Report e scritto possono essere fatti in appelli diversi; l'esame si registra quando
        entrambe le parti sono sufficienti.
\end{itemize}

\begin{esame}[come si scrive il report]
Le indicazioni del docente su come scrivere un buon report e la checklist prima della consegna sono
raccolte nel cap.~\ref{cap:lab-report}.
\end{esame}

\subsection{Uso di ChatGPT e strumenti simili}
È \textbf{permesso} usare l'AI per idee, bozze di codice e struttura del report. Restate però
responsabili di verificare e capire tutto ciò che consegnate: niente copia-incolla. Se la consegna
sembra interamente ``GPT-made'', senza ragionamento, il gruppo viene \textbf{chiamato per un orale}.
I testi dei lab aggiungono che saper trovare e correggere bug, anche nel codice generato dall'AI,
può essere valutato allo scritto, dove c'è una parte di comprensione del codice da fare senza AI.

\section{Esame scritto}
\begin{itemize}
  \item In aula, sulla piattaforma \textbf{Exam (Moodle)} con \textbf{LockDown Browser},
        \textbf{con il proprio portatile} (obbligatorio). Aggiorna LockDown Browser all'ultima versione.
  \item Copre \textbf{teoria e codice}: non si chiede di scrivere codice da zero, ma di capire se
        un codice funziona, perché no e cosa correggere.
  \item \textbf{Durata: 100 minuti} (in genere si finisce in 70--80).
  \item \textbf{Totale: 32 punti} secondo la slide (33 a voce, vedi riquadro iniziale).
\end{itemize}

\textbf{Struttura (descritta in aula):}
\begin{enumerate}
  \item domande a \textbf{scelta multipla}, due delle quali \textbf{senza penalità};
  \item \textbf{quattro gruppi di vero/falso}: risposta giusta = punto, risposta sbagliata = \textbf{piccola penalità};
  \item \textbf{due domande aperte} di teoria;
  \item \textbf{tre esercizi aperti} di difficoltà crescente sulla pipeline di ML e sul codice.
\end{enumerate}

\begin{verifica}
Dalla registrazione non è chiaro se le domande a scelta multipla siano \textbf{due in tutto} o se solo
due siano senza penalità; la slide dice soltanto ``\en{some with penalties for wrong answers}''.
\end{verifica}

\begin{esame}[cosa conta davvero]
Il docente ha ripetuto più volte che \textbf{non chiederà di ricordare le formule} delle slide:
conta capire \emph{perché} funzionano (es.\ perché il momentum stabilizza la traiettoria,
perché la sigmoid causa il \en{vanishing gradient}). Le formule nel manuale servono a
ragionare, non a essere memorizzate. Esempi di domande nei cap.~\ref{cap:quiz} e~\ref{cap:esercizi}.
\end{esame}

\section{Voto finale}
\[
\text{Voto} = 0{,}6 \cdot \text{Scritto} + 0{,}4 \cdot \overline{\text{Report}} + \text{punti aggiuntivi}
\]
\begin{itemize}
  \item Si arrotonda all'intero più vicino \textbf{solo alla fine} ($\geq 0{,}5$ per eccesso).
  \item Registrazione solo se \textbf{scritto $\geq 18$} \textbf{e ogni report $\geq 18$}.
  \item Punti aggiuntivi: +1 del Lab 1 consegnato entro il 1° dicembre ed eventuali altre attività proposte.
  \item \textbf{Orale facoltativo} se il voto è $\geq 26$, oppure su richiesta del docente in caso
        di dubbi sullo scritto o sul report.
\end{itemize}

\begin{esempio}[calcolo, dalla slide]
$24 \cdot 0{,}6 + 28 \cdot 0{,}4 + 1 = 14{,}4 + 11{,}2 + 1 = 26{,}6 \Rightarrow 27$\\
$20 \cdot 0{,}6 + 28 \cdot 0{,}4 + 1 = 12 + 11{,}2 + 1 = 24{,}2 \Rightarrow 24$
\end{esempio}

\begin{attenzione}[errore nella slide]
La slide riporta $23{,}2 \Rightarrow 23$ per il secondo esempio; il conto corretto è
$24{,}2 \Rightarrow 24$.
\end{attenzione}

\section{Mappa lezioni $\leftrightarrow$ slide}
\begin{center}
\small
\begin{tabularx}{\textwidth}{l X l}
\toprule
\textbf{Lezione} & \textbf{Argomenti} & \textbf{Slide}\\
\midrule
\rif{2026-09-21}{21/09} & Presentazione; perché AI in cybersecurity; AI/ML/NN/DL; storia; shallow vs deep & 0, 1 (pp.\ 1--25)\\
\rif{2026-09-23}{23/09} & Terminologia; paradigmi di apprendimento; uso dei modelli; AI \& cybersecurity & 1 (pp.\ 26--54)\\
\rif{2026-09-28}{28/09} & Perceptron, bias, attivazioni, stacking, non linearità, approssimazione universale, logit/softmax & 2 (pp.\ 1--18)\\
\rif{2026-09-30}{30/09} & Training: epoch, mini-batch, learning rate, GD/SGD/momentum/Adam, backprop, loss, metriche, over/underfitting & 2 (pp.\ 19--51)\\
\rif{2026-10-05}{05/10} & Data handling: pulizia, encoding, scaling, split, sbilanciamento, leakage; progettare una FFNN; tensori & 3 (pp.\ 1--22)\\
\rif{2026-10-06}{06/10} & FFNN in PyTorch: classe, training loop, mini-batch, early stopping, dropout, weight decay, GPU & 3 (pp.\ 23--33)\\
\rif{2026-10-12}{12/10} & RNN: hidden state, architetture, BPTT, gradienti, LSTM, dati 3D & 3 (pp.\ 34--71)\\
\rif{2026-10-14}{14/10} & RNN in PyTorch (padding); CNN; GNN (teoria) & 3 (pp.\ 72--123)\\
\rif{2026-10-19}{19/10} & GNN in PyTorch Geometric; feature categoriche ed embedding appresi & 3 (pp.\ 124--136), 4 (pp.\ 1--22)\\
\rif{2026-10-20}{20/10} & Embedding supervisionati e self-supervised; dropout e batch norm nel dettaglio; linee guida report & 4 (pp.\ 23--51), 2a\\
\rif{2026-11-02}{02/11} & Anomaly detection: definizioni, tipi, sfide; one-class SVM; clustering e k-means & 5 (pp.\ 1--33)\\
\rif{2026-11-09}{09/11} & Notebook di anomaly detection su KDD99: preprocessing, PCA/t-SNE, OC-SVM, LOF, k-means, DBSCAN, Isolation Forest, autoencoder & 5 (pp.\ 33--57), 6\\
\bottomrule
\end{tabularx}
\end{center}
